\documentclass[9pt]{extarticle}
\usepackage[margin=1.5cm]{geometry}
\usepackage[utf8]{inputenc}
\usepackage[spanish]{babel}
\usepackage{hyperref}
\usepackage{booktabs}
\usepackage{array}
\usepackage{titlesec}
\usepackage{enumitem}
\setlist{nosep, leftmargin=*, itemsep=1pt, topsep=1pt}
\titlespacing*{\section}{0pt}{6pt}{3pt}
\titlespacing*{\subsection}{0pt}{4pt}{2pt}
\setlength{\parskip}{2pt}
\setlength{\parindent}{0pt}
\pagestyle{empty}

\title{\vspace{-1.5cm}\large CC3092 -- Laboratorio \#1: Entrenamiento de Redes Neuronales (MLP)}
\author{Anthony Lou Schwank -- 23410}
\date{}

\begin{document}
\maketitle
\vspace{-0.9cm}

\textbf{Repositorio:} \url{https://github.com/anthonylouschwank/Lab1-DeepLearning}

\section*{1. Dataset y preparación}
Se usó \texttt{California Housing Prices} vía \texttt{sklearn.datasets.fetch\_california\_housing}: 20{,}640 observaciones, 8 features numéricas continuas (\texttt{MedInc}, \texttt{HouseAge}, \texttt{AveRooms}, \texttt{AveBedrms}, \texttt{Population}, \texttt{AveOccup}, \texttt{Latitude}, \texttt{Longitude}) y target \texttt{median\_house\_value}. Sin nulos ni duplicados. Se detectaron outliers reales (no errores) en \texttt{AveBedrms}, \texttt{Population}, \texttt{AveOccup} y \texttt{MedInc} vía regla 1.5$\times$IQR, y el target está \emph{censurado}: 4.68\% de las observaciones tope\'an en el valor máximo (\$500{,}001). No se eliminaron (son válidos y reducirían la muestra); se mitigó su efecto escalando features. No hay variables categóricas (la versión sklearn no incluye \texttt{ocean\_proximity}), por lo que no se requirió codificación. Se escalaron las 8 features con \texttt{StandardScaler} (ajustado solo sobre train) por las diferencias de escala entre variables (p.ej. \texttt{Population} en miles vs \texttt{Latitude} en decenas). División 70/15/15 (train/val/test): test separado primero y usado una única vez al final.

\section*{2. Capas y optimizadores de PyTorch (resumen)}
\textbf{\texttt{nn.Linear(in,out)}}: transformación afín $y=xW^T+b$, bloque básico fully-connected. \textbf{Activaciones}: \texttt{ReLU} ($\max(0,x)$, barata pero sufre \emph{dying ReLU}); \texttt{LeakyReLU} (permite gradiente pequeño para $x<0$, mitiga lo anterior); \texttt{Tanh} (acotada en $(-1,1)$, satura en redes profundas). \textbf{\texttt{nn.Dropout(p)}}: apaga aleatoriamente una fracción $p$ de neuronas en entrenamiento (regularización); se desactiva en \texttt{eval()}. \textbf{\texttt{nn.BatchNorm1d}}: normaliza activaciones por mini-batch (media 0, var 1) + transformación afín aprendible; estabiliza y acelera el entrenamiento. \textbf{Pérdidas}: \texttt{MSELoss} (sensible a outliers, penaliza cuadráticamente), \texttt{L1Loss} (robusta a outliers, gradiente constante), \texttt{SmoothL1Loss} (híbrida, MSE cerca de 0 y L1 lejos).

\textbf{Optimizadores}: todos comparten \texttt{lr} (tamaño de paso; muy alto diverge, muy bajo estanca) y \texttt{weight\_decay} (L2, penaliza pesos grandes). \texttt{SGD} actualiza en dirección del gradiente del mini-batch, opcionalmente con \texttt{momentum}; es el más simple pero requiere más ajuste manual. \texttt{Adam} combina momentum con tasa adaptativa por parámetro (media móvil de gradiente y de su cuadrado); robusto y rápido con poco ajuste, es el más usado por defecto. \texttt{RMSprop} adapta el \texttt{lr} dividiendo por la raíz de una media móvil del gradiente al cuadrado; predecesor directo de Adam, útil en problemas no estacionarios.

\section*{3. Resultados: 14 iteraciones (ordenadas por MSE de validación)}
Baseline: \texttt{[64,32]}, ReLU, Adam (lr=0.001), batch=32, epochs=40, sin regularización. Búsqueda tipo grid dirigido: cada iteración cambia 1--2 variables respecto al baseline.

\begin{center}
\begin{tabular}{@{}clllllccc@{}}
\toprule
\textbf{\#} & \textbf{Cambio} & \textbf{Arq.} & \textbf{Activ.} & \textbf{Optim/LR} & \textbf{Batch/Ep.} & \textbf{MSE} & \textbf{MAE} & \textbf{RMSE} \\
\midrule
10 & Batch pequeño & [64,32] & relu & adam/0.001 & 8/40 & \textbf{0.2680} & 0.3441 & 0.5177 \\
9  & Batch grande + más epochs & [64,32] & relu & adam/0.001 & 128/80 & 0.2696 & 0.3509 & 0.5193 \\
8  & LR alto (Adam) & [64,32] & relu & adam/0.01 & 32/40 & 0.2792 & 0.3585 & 0.5284 \\
13 & Regularización L1 & [64,32] & relu & adam/0.001 & 32/40 & 0.2797 & 0.3513 & 0.5289 \\
2  & Arq. más profunda & [128,64,32] & relu & adam/0.001 & 32/40 & 0.2822 & 0.3507 & 0.5312 \\
6  & SGD + momentum & [64,32] & relu & sgd/0.01 & 32/40 & 0.2826 & 0.3578 & 0.5316 \\
4  & Activación LeakyReLU & [64,32] & leaky\_relu & adam/0.001 & 32/40 & 0.2832 & 0.3548 & 0.5321 \\
7  & Optimizador RMSprop & [64,32] & relu & rmsprop/0.001 & 32/40 & 0.2834 & 0.3636 & 0.5324 \\
5  & Activación Tanh & [64,32] & tanh & adam/0.001 & 32/40 & 0.2884 & 0.3604 & 0.5371 \\
1  & Baseline & [64,32] & relu & adam/0.001 & 32/40 & 0.2893 & 0.3563 & 0.5379 \\
12 & Regularización L2 & [64,32] & relu & adam/0.001 & 32/40 & 0.2943 & 0.3624 & 0.5425 \\
3  & Arq. pequeña (capacidad limitada) & [16] & relu & adam/0.001 & 32/40 & 0.3186 & 0.3887 & 0.5644 \\
11 & Dropout + BatchNorm & [64,32] & relu & adam/0.001 & 32/40 & 0.3625 & 0.4074 & 0.6021 \\
14 & Combo: profunda+LeakyReLU+Dropout+L2 & [128,64,32] & leaky\_relu & adam/0.001 & 32/60 & 0.4251 & 0.4214 & 0.6520 \\
\bottomrule
\end{tabular}
\end{center}
\vspace{2pt}
\textbf{Modelo final (\#10) sobre test (evaluación única):} MSE = 0.2761, MAE = 0.3523, RMSE = 0.5254 -- consistente con validación, sin señales de que la selección haya sobreajustado a val.

\section*{4. Análisis de curvas de pérdida}
\textbf{\#1 (Baseline):} train y val bajan juntos y cierran cerca (train=0.266, val=0.289); leve subida de val tras el epoch 36, señal temprana y suave de overfitting. \textbf{\#3 (arq. pequeña):} train y val se estancan juntos en $\sim$0.32, gap$\approx$0: \emph{underfitting} claro, el modelo carece de capacidad. \textbf{\#10 (mejor):} train (medido en mini-batches ruidosos de tamaño 8) termina por encima del val (0.305 vs 0.268); es un artefacto de medición, no mal ajuste -- el modelo generaliza mejor que todos los demás. \textbf{\#14 (regularización fuerte):} mejor val en epoch 31 (0.294) pero se sigue entrenando hasta el epoch 59 y val sube a 0.425 mientras train baja; \emph{overfitting tardío} pese a la regularización -- \emph{early stopping} habría sido más efectivo que agregar más regularización.

\section*{5. Discusión}
\textbf{Mayor impacto positivo/negativo:} el cambio de \texttt{batch\_size} (chico, \#10, o grande+más epochs, \#9) dio la mayor mejora sobre el baseline. El mayor impacto negativo fue combinar dropout+BatchNorm (\#11) y sobre todo el combo con L2 y más epochs (\#14): el modelo pasó de generalizar bien a underfitting/overfitting tardío.

\textbf{Overfitting/underfitting:} sí, ver \#3 (underfitting por poca capacidad) y \#14 (overfitting tardío por exceder epochs necesarios). Se identificó comparando las curvas train vs val por epoch; la mitigación recomendada es \emph{early stopping} más que agregar regularización.

\textbf{¿Ayudó la regularización?} No, en ningún caso mejoró sobre el baseline. L1 (\#13, $\lambda=10^{-4}$) fue casi neutro por ser muy débil; L2 (\#12) y dropout+BatchNorm (\#11, \#14) empeoraron las métricas. Esto ocurre porque el baseline ya generaliza razonablemente bien (gap train/val pequeño), así que añadir restricciones de capacidad solo introduce underfitting en vez de corregir un overfitting que apenas existía.

\textbf{Batch size y epochs:} batch chico (8) genera más actualizaciones por epoch y una pérdida de entrenamiento más ruidosa, pero terminó generalizando mejor; batch grande (128) da un entrenamiento más estable/suave pero necesitó el doble de epochs (80) para acercarse al mismo desempeño, ya que hay menos actualizaciones por epoch.

\textbf{MSE vs MAE vs RMSE:} MSE/RMSE penalizan más los errores grandes (sensibles a outliers/target censurado detectado en el EDA); MAE da la magnitud típica del error de forma robusta. En el modelo final RMSE (0.525) $>$ MAE (0.352), lo que indica presencia de errores puntuales más grandes que el promedio -- coherente con la censura del target en \$500,001.

\textbf{Modelo de producción:} arquitectura \texttt{[64,32]} con Adam (lr=0.001), sin dropout/L2 agresivos (no ayudaron aquí), usando \emph{early stopping} sobre val en vez de regularización fuerte, y batch pequeño u optando por más epochs con batch grande. Para seguir optimizando usaría \textbf{búsqueda bayesiana} (p.ej. Optuna), más eficiente que grid/random para explorar simultáneamente batch, lr y arquitectura con un presupuesto de iteraciones limitado.

\section*{6. Conclusiones}
El MLP con arquitectura moderada (\texttt{[64,32]}) y Adam ya generaliza bien en este dataset sin necesitar regularización explícita; el principal margen de mejora vino de ajustar batch size/epochs, no de agregar restricciones de capacidad. El target censurado y los outliers reales (no removidos) explican por qué RMSE $>$ MAE consistentemente. El proceso respetó la separación estricta train/val/test, evaluando el conjunto de test una única vez sobre el modelo ya seleccionado.

\end{document}
